(1) 实验室坐标系中一段长 $L$ 的透明管内介质折射率 $n>1$，沿 $+x$ 以匀速 $u$ 流动且 $0<u\ll c$。光在介质静止系中的相速度为 $c/n$。以狭义相对论速度合成求实验室中顺流及逆流光的有符号相速度。
(2) 展开至 $u/c$ 的一阶，推导Fresnel拖曳系数 $1-1/n^2$。
(3) 假设两个光束在实验室固定的管长 $L$ 内反向传播并重新干涉，求一阶传播时间差 $\Delta t$ 与真空波长 $\lambda_0$ 下的条纹位移 $\Delta N$。
(4) 检验 $n\to1$ 与 $u\to0$ 的极限，说明色散介质的相、群速度问题及实验端点运动会如何修正简单模型。
